\subsection{\texorpdfstring{The complexes \(\mathsf{K}(u),\mathsf{K}.(u,\mathsf{C}),\mathsf{K}^{\prime}(u,\mathsf{C})\)}{The complexes \textbackslash mathsf\{K\}(u),\textbackslash mathsf\{K\}.(u,\textbackslash mathsf\{C\}),\textbackslash mathsf\{K\}\^{}\{\textbackslash prime\}(u,\textbackslash mathsf\{C\})}}

Let A be a ring, L an A-module, \(u: \mathbf{L} \to \mathbf{A}\) a linear form, and \(\Lambda(\mathbf{L})\) the exterior algebra of the A-module L. For \(x \in \Lambda(\mathbf{L})\)  let \(d_u(x)\) denote the interior product \(x \sqsubseteq u\) (III, p.~161, example). According to loc. cit., p.~162, formula (60), we have

\[
d _ {u} \left(e _ {1} \wedge \dots \wedge e _ {n}\right) = \sum_ {i = 1} ^ {n} (- 1) ^ {i + 1} u \left(e _ {i}\right) e _ {1} \wedge \dots \wedge e _ {i - 1} \wedge e _ {i + 1} \wedge \dots \wedge e _ {n}\tag{1}
\]

for \(e_1, \ldots, e_n\) in L. According to III, p.~164 and 165, the map \(d_u: \symbf{\Lambda}(\mathbf{L}) \to \symbf{\Lambda}(\mathbf{L})\) is an antiderivation of degree \((-1)\) and of square zero. It is the unique antiderivation of the A-algebra \(\symbf{\Lambda}(\mathbf{L})\) which extends \(u: \symbf{\Lambda}^1(\mathbf{L}) \to \symbf{\Lambda}^0(\mathbf{L})\) .

DEFINITION 1. --- The complex \((\symbf{\Lambda}(\mathbf{L}), d_u)\) is denoted \(\mathbf{K}^{\mathbf{A}}(u)\) or \(\mathbf{K}(u)\) .

Be careful that \(\mathbf{K}_{n}(u)=\symbf{\Lambda}^{n}(\mathbf{L})=\mathbf{K}^{-n}(u)\) It is clear that \(\mathbf{K}(u)\) is zero on the right and that \(\mathrm{H}_{0}(\mathbf{K}(u))=\mathrm{Coker}(u)=\mathrm{A}/\mathfrak{q}\) where \(\mathfrak{q}\) is the ideal \(u(\mathbf{L})\) of A.

For every complex C of A-modules, we set

\[
\mathbf {K} _ {\cdot} ^ {\mathrm{A}} (u, \mathrm{C}) = \mathrm{C} \otimes_ {\mathrm{A}} \mathbf {K} ^ {\mathrm{A}} (u), \quad \mathbf {K} _ {\cdot} ^ {\cdot} (u, \mathrm{C}) = \operatorname{Homgr} _ {\mathrm{A}} (\mathbf {K} ^ {\mathrm{A}} (u), \mathrm{C}).
\]

\[
\mathrm{H} ^ {\mathsf {A}} (u, \mathrm{C}) = \mathrm{H} (\mathrm{C} \otimes_ {\mathsf {A}} \mathbf {K} ^ {\mathsf {A}} (u)) \quad \mathrm{H} _ {\mathsf {A}} ^ {\bullet} (u, \mathrm{C}) = \mathrm{H} (\mathrm{Homgr} _ {\mathsf {A}} (\mathbf {K} ^ {\mathsf {A}} (u), \mathrm{C})) ,
\]

\[
\mathrm{H} _ {r} ^ {\mathrm{A}} (u, \mathrm{C}) = \mathrm{H} _ {r} (\mathrm{C} \otimes_ {\mathrm{A}} \mathbf {K} ^ {\mathrm{A}} (u)), \quad \mathrm{H} _ {\mathrm{A}} ^ {r} (u, \mathrm{C}) = \mathrm{H} ^ {r} (\operatorname{Homgr} _ {\mathrm{A}} (\mathbf {K} ^ {\mathrm{A}} (u), \mathrm{C})) .
\]

We therefore have canonical homomorphisms of A-modules (X, p.~62 and p.~82)

\[
\gamma_ {0}: \mathrm{H} _ {0} (\mathbf {C}) \otimes_ {\mathbf {A}} \mathrm{A} / \mathfrak {q} \rightarrow \mathrm{H} _ {0} ^ {\mathbf {A}} (u, \mathbf {C}),
\]

\[
\lambda^ {0}: \mathrm{H} _ {\mathrm{A}} ^ {0} (u, \mathrm{C}) \rightarrow \operatorname{Hom} _ {\mathrm{A}} (\mathrm{A} / \mathfrak {q}, \mathrm{H} ^ {0} (\mathrm{C})).
\]

Lemma 1. --- If the complex C is zero on the right (resp. on the left), then \(\mathbf{K}^{\mathbf{A}}(u,\mathbf{C})\) (resp. \(\mathbf{K}_{\mathbf{A}}^{*}(u,c)\) ) is zero on the right (resp. on the left), and \(\gamma_{0}\) (resp. \(\lambda^{0}\) ) is bijective.

This follows from X, p.~62, prop. 1 and p.~82, prop. 1.

PROPOSITION 1. --- Let \(x \in L\) ; denote by \(\mathbf{R}_{x}: y \mapsto x \wedge y\) the left multiplication by \(x\) in the algebra \(\boldsymbol{\Lambda}(\mathbf{L})\) . Then \(d_{u} \circ \mathbf{R}_{x} + \mathbf{R}_{x} \circ d_{u} = u(x)\) . \(1_{\boldsymbol{\Lambda}(\mathbf{L})} = u(x)_{\boldsymbol{\Lambda}(\mathbf{L})}\) .

Indeed, \((d_u \circ R_x + R_x \circ d_u)(y) = d_u(x \wedge y) + x \wedge d_u(y)\) ; since \(d_u\) is a derivation, \(d_u(x \wedge y) + x \wedge d_u(y) = d_u(x) \wedge y = u(x).y\) .

COROLLARY 1. --- If u is surjective, \(\mathbf{K}(u)\) is homotopic to zero (X, p.~34) as is \(\mathbf{K}_{\cdot}^{\mathrm{A}}(u,\mathrm{C})\) and \(\mathbf{K}_{\cdot}^{\cdot}(\mathbf{u},\mathbf{C})\) for every complex C.

Indeed, there exists \(x \in \mathbf{L}\) such that \(u(x) = 1\) . Then \(\mathbf{K}(u)\) is homotopic to zero by prop. 1, hence also \(\mathbf{K}^{\mathrm{A}}(u, \mathbf{C})\) (X, p.~64, prop. 3) and \(\mathbf{K}_{\mathrm{A}}^{*}(u, \mathbf{C})\) (X, p.~83, prop. 3).

COROLLARY 2. --- Let C be a complex, Ann (C) its annihilator. Then q + Ann (C) annihilates H \(^{A}\) (u, C) and H \(^{*}_{A}\) (u, C).

For every \(\lambda \in q\) , the homothety \(\lambda_{\mathbf{K}(u)}\) is homotopic to zero by the proposition, hence also \(1_{C} \otimes \lambda_{\mathbf{K}(u)}\) and Homgr \((\lambda_{\mathbf{K}(u)}, 1_{C})\) by X, p.~64, prop. 3 and X, p.~83, prop. 3; it follows that \(\lambda\) annihilates H.(u, C) and H'(u, C). If \(\lambda \in \text{Ann}(C)\) , then \(1_{\mathbf{K}(u)} \otimes \lambda_{C}\) and Homgr \((1_{\mathbf{K}(u)}, \lambda_{C})\) are zero.

Suppose L is projective (resp. \(\mathbf{K}(u)\) acyclic in degrees \(>0\) ). Then the complex \(\symbf{\Lambda}(\mathbf{L})\) is projective by III, p.~87, cor. 2 (resp. is a resolution of A/q); by X, p.~102 (resp. p.~100), we therefore have, for every A-module M, homomorphisms

\begin{enumerate}
\def\labelenumi{(\arabic{enumi})}
\setcounter{enumi}{1}
\tightlist
\item
\end{enumerate}

\[
\mathrm{H} _ {r} ^ {\mathbf {A}} (u, \mathrm{M}) \rightarrow \operatorname{Tor} _ {r} ^ {\mathbf {A}} (\mathrm{A} / \mathrm{q}, \mathrm{M}), \quad \operatorname{Ext} _ {\mathbf {A}} ^ {r} (\mathrm{A} / \mathrm{q}, \mathrm{M}) \rightarrow \mathrm{H} _ {\mathbf {A}} ^ {r} (u, \mathrm{M})\tag{3}
\]

\[
\operatorname{Tor} _ {r} ^ {\mathrm{A}} (\mathrm{A} / \mathrm{q}, \mathrm{M}) \rightarrow \mathrm{H} _ {r} ^ {\mathrm{A}} (u, \mathrm{M}), \quad \mathrm{H} ^ {r} (u, \mathrm{M}) \rightarrow \operatorname{Ext} _ {\mathrm{A}} ^ {r} (\mathrm{A} / \mathrm{q}, \mathrm{M}).
\]

If L is projective and \(\mathbf{K}(u)\) acyclic in degrees \(>0\) , the homomorphisms (2) and (3) above are bijective and inverses of each other (X, p.~102, prop. 1).

PROPOSITION 2. --- Let \((\mathbf{L}_{i})_{i\in I}\) be a family of A-modules, where the set I is finite and totally ordered. Let u be a linear form on \(\bigoplus_{i\in I}\mathbf{L}_{i}\) , \(u_{i}\) its restriction to \(\mathbf{L}_{i}\) .

The canonical isomorphism of A-algebras (III, p.~84)

\[
g: \bigotimes_{\substack{i\in I}}\symbf{\Lambda}(\mathrm{L}_{i})\to \symbf{\Lambda}(\bigoplus_{i\in I}\mathrm{L}_{i})
\]

is an isomorphism of the complex \(\otimes \mathbf{K}(u_i)\) (X, p.~63) on the complex \(\mathbf{K}(u)\) .

Indeed, according to X, p.~64, remark 4, the differential D of the complex \(\otimes_{i\in I}\mathbf{K}(u_i)\) is an antiderivation; the antiderivations \(d_u\) and \(g\circ D\circ g^{-1}\) of \(\boldsymbol{\Lambda}(\oplus L_i)\) coincide on \(\oplus L_i\) with the map \(x\mapsto u(x)\) . 1 of \(\oplus L_i\) to \(\boldsymbol{\Lambda}(\oplus L_i)\) , hence they are equal (III, p.~128, cor.).

Let C and C' be two complexes of A-modules. We have (X, p.~63 and p.~99) canonical isomorphisms of complexes

\[
\mathbf {C} \otimes_ {\mathbf {A}} \left(\mathbf {C} ^ {\prime} \otimes_ {\mathbf {A}} \mathbf {K} (u)\right)\rightarrow \left(\mathbf {C} \otimes_ {\mathbf {A}} \mathbf {C} ^ {\prime}\right) \otimes_ {\mathbf {A}} \mathbf {K} (u)
\]

\[
\operatorname{Homgr} _ {\mathbf {A}} \left(\mathrm{C} ^ {\prime}, \operatorname{Homgr} _ {\mathbf {A}} (\mathbf {K} (u), \mathrm{C})\right)\rightarrow \operatorname{Homgr} _ {\mathbf {A}} \left(\mathrm{C} ^ {\prime} \otimes_ {\mathbf {A}} \mathbf {K} (u), \mathrm{C}\right),
\]

that is, isomorphisms

\begin{enumerate}
\def\labelenumi{(\arabic{enumi})}
\setcounter{enumi}{3}
\tightlist
\item
\end{enumerate}

\[
\mathbf {C} \otimes_ {\mathbf {A}} \mathbf {K} ^ {\mathrm{A}} (u, \mathbf {C} ^ {\prime}) \rightarrow \mathbf {K} ^ {\mathrm{A}} (u, \mathbf {C} \otimes_ {\mathbf {A}} \mathbf {C} ^ {\prime})\tag{5}
\]

\[
\operatorname{Homgr} _ {\mathbf {A}} \left(\mathrm{C} ^ {\prime}, \mathbb {K} _ {\mathbf {A}} ^ {*} (u, \mathrm{C})\right)\rightarrow \operatorname{Homgr} _ {\mathbf {A}} \left(\mathbb {K} _ {\mathbf {A}} ^ {*} (u, \mathrm{C} ^ {\prime}), \mathrm{C}\right).
\]

In (4) and (5), take \(\mathbf{C}^{\prime} = \mathbf{K}(u^{\prime})\) , where \(u^{\prime}: L^{\prime} \to A\) is a linear form on an A-module \(L^{\prime}\) and note that \(\mathbf{K}^{\mathrm{A}}(u, \mathbf{K}(u^{\prime}))\) which is equal by definition to \(\mathbf{K}(u^{\prime}) \otimes_{\mathrm{A}} \mathbf{K}(u)\) identifies, according to prop. 2, with \(\mathbf{K}(u^{\prime} \oplus u)\) where \(u^{\prime} \oplus u: L^{\prime} \oplus L \to A\) is the linear form \((x^{\prime}, x) \mapsto u^{\prime}(x^{\prime}) + u(x)\) We then obtain isomorphisms of complexes

\begin{enumerate}
\def\labelenumi{(\arabic{enumi})}
\setcounter{enumi}{5}
\tightlist
\item
\end{enumerate}

\[
\mathbf {K} _ {\bullet} ^ {\mathrm{A}} (u ^ {\prime} \oplus u, \mathrm{C}) \rightarrow \mathbf {K} _ {\bullet} ^ {\mathrm{A}} (u, \mathbf {K} _ {\bullet} ^ {\mathrm{A}} (u ^ {\prime}, \mathrm{C}))\tag{7}
\]

\[
\mathbf {K} _ {\mathrm{A}} ^ {\bullet} (u ^ {\prime}, \mathbf {K} _ {\mathrm{A}} ^ {\bullet} (u, \mathrm{C})) \rightarrow \mathbf {K} _ {\mathrm{A}} ^ {\bullet} (u ^ {\prime} \oplus u, \mathrm{C}).
\]

By passing to homology, we deduce isomorphisms of A-modules

\[
\begin{array}{c c} \mathrm{H} _ {r} ^ {\mathbf {A}} (u ^ {\prime} \oplus u, \mathrm{C}) \to \mathrm{H} _ {r} ^ {\mathbf {A}} (u, \mathbf {K} _ {\bullet} ^ {\mathbf {A}} (u ^ {\prime}, \mathrm{C}))  , & r \in \mathbf {Z}  , \\ \mathrm{H} _ {\mathbf {A}} ^ {r} (u ^ {\prime}, \mathbf {K} _ {\mathbf {A}} ^ {\bullet} (u, \mathrm{C})) \to \mathrm{H} _ {\mathbf {A}} ^ {r} (u ^ {\prime} \oplus u, \mathrm{C})  , & r \in \mathbf {Z}  . \end{array}
\]

Finally, note that the homomorphism deduced from the product in the algebra \(\Lambda(L)\)

\[
m: \mathbf {K} ^ {\mathrm{A}} (u) \otimes_ {\mathrm{A}} \mathbf {K} ^ {\mathrm{A}} (u) \rightarrow \mathbf {K} ^ {\mathrm{A}} (u)
\]

is a morphism of complexes (since \(d_{u}\) is an antiderivation). Assuming L is free of rank \(n\) and composing with the morphism of complexes \(\mathbf{K}^{\mathrm{A}}(u)\to\symbf{\Lambda}^{n}\mathrm{L}(-n)\) which is the identity in degree \(n\) from this one deduces a morphism of complexes

\[
\chi : \mathbf {K} ^ {\mathrm{A}} (u) \otimes_ {\mathrm{A}} \mathbf {K} ^ {\mathrm{A}} (u) \rightarrow \boldsymbol {\Lambda} ^ {n} \mathrm{L} (- n);
\]

to this morphism there corresponds canonically, according to X, p.~99, prop. 12, a morphism of complexes

\[
\varphi : \mathbf {K} ^ {\mathrm{A}} (u) \rightarrow \operatorname{Homgr} _ {\mathrm{A}} \left(\mathbf {K} ^ {\mathrm{A}} (u), \boldsymbol {\Lambda} ^ {n} \mathrm{L} (- n)\right)
\]

which is bijective (III, p.~87, formula (20)). For every complex C, one deduces from this a composed isomorphism

\[
\begin{array}{r l} \mathbf {K} _ {\bullet} ^ {\mathrm{A}} (u, \mathrm{C}) & = \mathrm{C} \otimes_ {\mathrm{A}} \mathbf {K} ^ {\mathrm{A}} (u) \xrightarrow {1 \otimes_ {\Phi}} \mathrm{C} \otimes \operatorname{Homgr} _ {\mathrm{A}} \left(\mathbf {K} ^ {\mathrm{A}} (u), \boldsymbol {\Lambda} ^ {n} \mathrm{L} (- n)\right) \to \\ & \to \operatorname{Homgr} _ {\mathrm{A}} \left(\mathbf {K} ^ {\mathrm{A}} (u), \mathrm{C} \otimes_ {\mathrm{A}} \boldsymbol {\Lambda} ^ {n} \mathrm{L} (- n)\right) = \mathbf {K} _ {\mathrm{A}} ^ {\bullet} (u, \mathrm{C} \otimes_ {\mathrm{A}} \boldsymbol {\Lambda} ^ {n} \mathrm{L} (- n)). \end{array}
\]

By passing to homology, one therefore has canonical isomorphisms

\[
\mathrm{H} _ {r} ^ {\mathbf {A}} (u, \mathrm{C}) \rightarrow \mathrm{H} _ {\mathbf {A}} ^ {n - r} (u, \mathrm{C} \otimes_ {\mathbf {A}} \boldsymbol {\Lambda} ^ {n} \mathrm{L}), \quad r \in \mathbf {Z}.\tag{8}
\]

Remarks. --- 1) * The preceding remains valid when L is projective of rank n.~*

\begin{enumerate}
\def\labelenumi{\arabic{enumi})}
\setcounter{enumi}{1}
\tightlist
\item
  Since L is free of rank n, \(\Lambda^{n}\) L is isomorphic to A, we have non-canonical isomorphisms \(\mathrm{H}_{r}^{\mathrm{A}}(u,\mathrm{C})\to\mathrm{H}_{\mathrm{A}}^{n-r}(u,\mathrm{C})\) .
\end{enumerate}

